% Společné nastavení pro protokoly KI/PZS (pdfLaTeX, Overleaf výchozí kompilátor)
\usepackage[utf8]{inputenc}
\usepackage[T1]{fontenc}
\usepackage[czech]{babel}
\usepackage{lmodern}
\usepackage{amsmath}
\usepackage[a4paper,margin=2.5cm]{geometry}
\usepackage{microtype}
\usepackage{graphicx}
\usepackage{booktabs}
\usepackage{array}
\usepackage{tabularx}
\usepackage{xcolor}
\usepackage{enumitem}
\usepackage[font=small,labelfont=bf]{caption}
\usepackage{subcaption}
\usepackage{float}
\usepackage{tikz}
\usetikzlibrary{arrows.meta,positioning,shapes.geometric,calc}
\usepackage{pgfplots}
\pgfplotsset{compat=1.17}
\usepackage[hidelinks]{hyperref}

\setlist{itemsep=2pt,topsep=4pt}
\setlength{\parindent}{0pt}
\setlength{\parskip}{6pt}

% Záporné znaménko v textových tabulkách: \m0,5 -> −0,5
\newcommand{\m}{$-$}

% Červená značka pro místa, která je třeba před odevzdáním doplnit
\newcommand{\doplnit}[1]{\textcolor{red}{\textbf{[DOPLNIT: #1]}}}

% Obrázek, který se zobrazí, jakmile je soubor nahrán; do té doby červený zástupce
\newcommand{\obrazekNeboZastupce}[3]{%
  \IfFileExists{#1}{\includegraphics[width=#2]{#1}}{%
    \fbox{\parbox[c][4cm][c]{\dimexpr#2-2\fboxsep-2\fboxrule\relax}{\centering\doplnit{#3}}}}}

% Styl vývojových diagramů
\tikzset{
  krok/.style={draw, rounded corners=3pt, minimum width=3.3cm, minimum height=1.15cm,
               align=center, font=\small, line width=0.5pt, draw=black!60},
  hlavni/.style={krok, draw=blue!70!black, fill=blue!8, line width=0.9pt},
  rozhodnuti/.style={draw, diamond, aspect=2.4, align=center, font=\small,
               inner sep=1pt, line width=0.5pt, draw=black!60},
  sipka/.style={-{Stealth[length=2.6mm]}, line width=0.5pt, draw=black!60},
  popisek/.style={font=\footnotesize, text=black!70}
}
